% !TeX root = _main.tex
% ========================================
% はじめに
% ========================================

%1
\section{はじめに}
音は，視覚以外の情報伝達手段として，日常生活で幅広く利用されている．踏切音や通知音は，視覚情報を伴わなくとも，特定の出来事の発生を即座に伝える役割を担っている．また，映画やアニメーションなどの映像作品には，登場人物の心理を音によって示す表現が存在する．このように，音は視覚に依存せず，出来事や心理的状態を伝える情報として機能している．

音による情報伝達に関する研究は数多く存在するが，それらの多くは，単発的な音や限定的な状況を対象としている．一方で，物語という複数のイベントの連続から構成される作品形式において，音によって意図された出来事や感情の知覚がどのように成立するのかについては，十分に検討されていない．音のみを用いた物語を聴取する場合，聴者は個々の音が何を示しているのかを知覚するだけでなく，それらが物語の文脈上でどのような意味を持つのかを理解する必要がある．しかし，音のみで構成された物語において，こうした理解がどこまで行われるかは，十分に整理されていない．この点を踏まえると，音のみで構成された表現の連続が，聴取時にどのように知覚されるのかを明らかにする必要がある．

そこで本研究では，音による情報伝達を物語表現に適用する試みとして，音のみを用いた物語作品を制作した．本研究の目的は，音のみで構成された物語表現において，物語上の出来事や感情を表現するために制作者が行った音声設計の考え方を整理し，それらを組み合わせた1つの物語を制作することである．
